\begin{table}[htbp]
  \centering
  \small
  \caption{Términos y N-gramas Más Frecuentes en el Corpus de Conversaciones}
  \label{tab:v2_ngrams}
  \begin{tabular}{llrrlrr}
    \toprule
    Tipo & N-grama & Conteo & Tipo & N-grama & Conteo \\
    \midrule
    Unigrama  & \textit{no}             & 4.124 & Bigrama  & \textit{alguien sabe}     & 1.192 \\
    Unigrama  & \textit{si}             & 2.871 & Bigrama  & \textit{grupo whatsapp}   & 456   \\
    Unigrama  & \textit{alguien}        & 2.315 & Bigrama  & \textit{buenas tardes}    & 412   \\
    Unigrama  & \textit{grupo}          & 1.942 & Bigrama  & \textit{curso temporada}  & 388   \\
    Unigrama  & \textit{materia}        & 1.830 & Trigrama & \textit{alguien sabe si}  & 425   \\
    Unigrama  & \textit{paralelo}       & 1.420 & Trigrama & \textit{enlace del grupo} & 231   \\
    \bottomrule
  \end{tabular}
  \vspace{4pt}
  \begin{minipage}{0.92\textwidth}
    \footnotesize
    \textit{Nota.} Datos calculados a partir de los 27.334 registros de frecuencias almacenados en la tabla \texttt{ngrams\_frecuencias}. Las stopwords estándar del español fueron excluidas a excepción de los marcadores de negación sintáctica (\emph{no}, \emph{nunca}, \emph{tampoco}).
  \end{minipage}
\end{table}
